% Integration fragment. Renumber the provisional B.7-L labels in the main paper.
% AI-produced by GPT-6 Astra (OpenAI, ChatGPT Pro mode), 2026-11-16.

\subsection{Two source checks in LM7\_PUSH E.64}

\paragraph{B.7-L1 (Known, with the reduction proved): the selected sewing stabilizer.}
Let $S$ be the closed oriented surface of genus two and let $\gamma$ be
a fixed oriented nonseparating simple closed curve. All mapping classes
preserve the orientation of the surface. The subscript $\vec\gamma$
requires preservation of the orientation of the curve as well.
Let $t$ be its positive Dehn twist. For the compact genus-one surface
$S_1^2$ with its two boundary components fixed pointwise, write $d$ for
the first positive boundary twist. In the standard positive three-chain
presentation, there is an identification
\[
 \Gamma(S_1^2,\partial S_1^2)=B_4\times\langle d\rangle,
 \qquad \delta_1=d,\qquad \delta_2=z_4d^{-1},
 \qquad z_4=(\sigma_1\sigma_2\sigma_3)^4.
\]
Margalit--Winarski, Section 6, gives $\operatorname{Mod}(S_1^2)
\cong B_4\times\mathbf Z$, with the extra generator a boundary twist;
the positive three-chain relation identifies the second boundary twist
as displayed. The notation here fixes both boundary components
pointwise, and differs from a convention allowing their permutation.

Gluing an annulus between these boundaries sends both $\delta_i$ to $t$.
Cutting a representative of a mapping class preserving $\vec\gamma$
produces a preimage. Its ambiguity is the opposite pair of boundary
twists $\delta_1\delta_2^{-1}$. Therefore the discrete sewing map has
kernel generated by $(z_4,-2)$, with an immaterial inversion of this
generator. Put
\[
 E=(\widehat B_4\times\widehat{\mathbf Z})/
       \overline{\langle(z_4,-2)\rangle},\qquad t_E=[(1,1)].
\]
The closed subgroup by which we divide is central. There is a natural
continuous sewing map $E\to(\widehat\Gamma(S))_{\vec\gamma}$.
This map is an isomorphism, including its central inertia, as follows.

First, $a\mapsto t_E^a$ is an injection of $\widehat{\mathbf Z}$.
The continuous homomorphism
\[
 (b,a)\longmapsto\operatorname{exp}(b)+6a
      \quad:\quad\widehat B_4\times\widehat{\mathbf Z}
           \longrightarrow\widehat{\mathbf Z}
\]
vanishes on $(z_4,-2)$, since $\operatorname{exp}(z_4)=12$.
It sends $t_E^a$ to $6a$. Multiplication by $6$ on
$\widehat{\mathbf Z}=\prod_p\mathbf Z_p$ is injective.
Second, quotienting $E$ by this inertia gives
\[
 E/t_E^{\widehat{\mathbf Z}}
   =\widehat B_4/\overline{\langle z_4\rangle}
   \cong\widehat{P\Gamma}_{1,2}.
\]
The last identification follows from the same boundary presentation by
capping the two boundaries by two individually labelled punctures.
Completion commutes with this specified quotient: finite quotients on
both sides are exactly the finite quotients of $B_4$ killing $z_4$.

For the intrinsic stabilizer, the opened source is Boggi,
\emph{Automorphisms of profinite mapping class groups}, arXiv
2011.15075v6, Theorem 2.3(i). In the present one-curve case its second
exact sequence reads
\[
 1\longrightarrow t^{\widehat{\mathbf Z}}
 \longrightarrow(\check\Gamma(S))_{\vec\gamma}
 \longrightarrow P\check\Gamma_{1,2}\longrightarrow1.
\]
The two newly created punctures are fixed individually, precisely
because the curve is oriented. Theorem 2.3 states that the left-hand
group is the profinite completion of the twist group. Proposition
2.6(i) states, in its own notation,
\[
 \check\Gamma(S,\partial S)=\widehat\Gamma(S,\partial S)
       \quad\text{for }g(S)\leq2,
\]
and the same equality for the pure groups. The paragraph immediately
preceding Theorem 2.3 also gives the genus-at-most-two comparison for
punctured surfaces. Thus the checks may be replaced by hats in the
displayed exact sequence. This is the exact profinite input needed;
no assertion about all-genus congruence is used.

The sewing map between these two exact sequences is the identity on
$\widehat{\mathbf Z}$ and on $\widehat{P\Gamma}_{1,2}$.
For completeness, if an element of its kernel is given, its image in
the quotient is trivial; hence it is $t_E^a$, and the equality on
inertia forces $a=0$. Conversely, for an element of the target, lift
its quotient to $E$. The difference lies in target inertia and is
also in the image. The map is therefore bijective and, being a
continuous bijection between profinite groups, is an isomorphism.
This proves the first line of E.64.3 with $t_E^2=z_4$.

For its action formula, use the standard positive genus-two chain
$T_1,\ldots,T_5$, put $t=T_5$ and identify the indicated $B_4$ with
$\langle T_1,T_2,T_3\rangle$. The three-chain relation gives
$(T_1T_2T_3)^4=T_5^2$. The element
$z_5=(T_1T_2T_3T_4)^5$ is the hyperelliptic involution: its sphere
quotient is trivial (the complementary disc contains only one marked
point), whereas its homological action is $-I$. In particular it is
central. Consequently $f(z_5,T_5^2)=1$ for every commutator word $f$.
The standard GT formula therefore sends $T_5$ to $T_5^\lambda$ and
restricts on the first three generators to $\Phi_{4,F}$.
On $E$ this is exactly
\[
 [(b,d^a)]\longmapsto[(\Phi_{4,F}(b),d^{\lambda a})].
\]
It respects $(z_4,-2)$ since $\Phi_{4,F}(z_4)=z_4^\lambda$.
Here the descent of the standard braid action to the genus-two
hyperelliptic quotient is the separate presentation argument in
E.64.2. This verifies its claimed \emph{selected} sewing square.
It does not identify an arbitrary orbifold action with a restriction
of this action.

\paragraph{B.7-L2 (Known, with the signs and central lifts derived): E.64.13.}
Use exactly the braid generators $(a,b,c)$ and Ferrari basis
\[
 \xi_1=ca^{-1},\qquad
 \xi_2=aba c^{-1}b^{-1}a^{-1}
\]
of E.64. Let $\theta:\widehat B_3^{\rm mor}\to\widehat B_4$
be its specific subgroup map with $\theta(u)=ab$. For
$\sigma\in G_{\mathbf Q}$ assume $\chi(\sigma)=1$.
Let $A_\sigma$ denote its positive pure order-three frame and
$\Omega_{3,\sigma}$ the original triangle action, lifted as in
E.64.28. The source statements are Nakamura--Tsunogai--Yasuda
(2010), Proposition 4.11 (p.444) and Proposition 4.13 (p.445).
Item 4.12 is a subsection, not a proposition.

In the source's notation $a=\tau_1$, $b=\tau_2$, $c=\tau_3$.
Proposition 4.11 gives at multiplier one
\[
 u_\sigma=h_\sigma(a^2,b^2)^{-1}(ab)^{3\rho_K(\sigma)}.
\]
The $\lambda=1$ specialization of the source's (1.1) is
$f(x,y)=h(y,z)^{-1}h(x,y)$. Its conjugacy equation identifies
$h_\sigma(a^2,b^2)^{-1}$ with $A_\sigma$ modulo the three-strand
centre; both lie in the same free pure complement, so uniqueness
of the positive pure frame gives equality there. Equivalently,
one can compare $f^{-1}(ab)=h^{-1}(ab)h$, which is the order-three
frame equation in this convention. Thus
\[
 u_\sigma=A_\sigma(ab)^{3\rho_K(\sigma)}.
\]
Proposition 4.13 gives, on the three triangle generators,
\[
 \Phi_{4,\sigma}\theta
       =\operatorname{Inn}(u_\sigma)\theta\Omega_{3,\sigma}.
\]
These are equations in the full completed braid group. In the
source proof the initially possible power of its centre is killed
by the exponent sum $\widehat B_4\to\widehat{\mathbf Z}$; the
full twist has exponent $12$, and multiplication by $12$ is
injective on $\widehat{\mathbf Z}$. No residual central sign
is silently discarded.

Normalize by the frame and apply $\theta^{-1}$. The homomorphism
$\theta$ sends $u^{3\rho_K}$ to $(ab)^{3\rho_K}$, so
\[
 \begin{split}
 \Psi_{3,\sigma}
 &=\theta^{-1}\operatorname{Inn}(A_\sigma^{-1})
                      \Phi_{4,\sigma}\theta\\
 &=\operatorname{Inn}(u^{3\rho_K(\sigma)})\Omega_{3,\sigma},
 \end{split}
\]
and hence
\[
 \Psi_{3,\sigma}\Omega_{3,\sigma}^{-1}
     =\operatorname{Inn}(u^{3\rho_K(\sigma)}).
\]
On $\widehat K^{\rm ab}$ the conjugation action of $ab$ is
$R_+=UL=\left(\begin{smallmatrix}0&1\\-1&1\end{smallmatrix}\right)$,
with $R_+^3=-I$. This proves
\[
 \mathcal D_3(\sigma)=(-I)^{\rho_K(\sigma)}.
\]
The central element $(ab)^3$ of the first three-strand subgroup
is distinct from $z_4$. It acts by $-I$ on the Ferrari homology,
whereas $z_4$ acts trivially. This distinction explains the sign.

The parity can be compared to the referee's common radical without
any informal division by two. Definition 4.7 (p.443) fixes
\[
 K^{1/n}=3^{1/(2n)}4^{-1/n}\zeta_{12n}^{-5}.
\]
For $n=2$, with $r_0=(3/4)^{1/4}>0$, this is
$K^{1/2}=(r_0/\sqrt2)\zeta_{24}^{-5}$. The other factors and
$K$ itself lie in $\mathbf Q^{\rm ab}$. Therefore
\[
 (-1)^{\rho_K(\sigma)}=\frac{\sigma(r_0)}{r_0}
       \qquad(\sigma\in G_{\mathbf Q^{\rm ab}}).
\]
This is the same bit used in the repaired Galois part of E.65.

\paragraph{B.7-L3 (Undecided): the all-GT comparison is not supplied by those sources.}
Proposition 4.4 of the same paper constructs the original triangle
Tate action for every $F\in\widehat{\mathrm{GT}}$. Propositions
4.11 and 4.13, and Theorem 4.16, explicitly have the different
domain $G_{\mathbf Q}$. The introduction, p.432, says:
\begin{quote}
``Unknown is whether they hold for all elements of GT''.
\end{quote}
Here ``they'' refers to the power-conjugate braid formulas of
Proposition 4.13. Theorem 4.16 gives the matrix comparison on
Galois and cannot be extended by replacing its quantifier.

Likewise, Lochak--Nakamura--Schneps (2026), Theorem 1.1, gives
the full profinite mapping-class actions with at most one marked
point. Its Section 8 explicitly leaves the required preservation
of cut-surface images and Lego compatibility for further work.
Boggi's stabilizer theorem identifies the \emph{groups}; it does
not identify two independently constructed actions on them.
The exact missing assertion for A.4(e) is an all-GT computation
of the comparisons $C_3(F)B_{3,F}^{-1}$ and
$C_2(F)B_{4,F}^{-1}$, with the stated paths and normalizations.
Neither the selected sewing square nor the cited full-domain
triangle action makes these quotients trivial.

\paragraph{B.7-L4 (A convention correction, not a change to the bit).}
Nakamura--Schneps (2000), Proposition 2.4 and Definition 5.1,
give the abelianization of their harmonic parameter as
$g_\sigma^{\rm ab}=(xy)^{\rho_2(\sigma)}$, where $\rho_2$
uses the positive roots of $2$. With the E.64 dictionary
$g=\beta^{-1}$ and the E.65 definition of $\rho$ from
$\beta^{-1}$, this is $\rho=\rho_2$, not $-\rho_2$.
The minus sign in the referee report's Marin-convention remark
should therefore be removed after identifying its generators
with the source's generators. The positive shear $+8\rho_2$
in the proof of NTY Theorem 4.16 independently checks this
choice. Its reduction modulo two is unaffected by changing
sign, so the correction changes neither $\delta$ nor $\Delta$.
